\subsection{\problemblock{*Initial} Data Block}\label{sec:block-initial}

The initial conditions, or starting point, of a trajectory are given in the
\problemblock{*initial} data block. They define the initial vehicle state: position,
velocity, and mass. They may also define starting time, ground range, and path
length.

Initial conditions can be provided directly by specifying position, velocity, and
mass, or indirectly by taking them from another trajectory. The second method is
used when a vehicle releases or drops an object or when a missile stages, allowing
all resulting trajectories to be computed together.

\taossubhead{Specific Initial Conditions}

When specific conditions are entered, the block begins with a coordinate-system
keyword followed by variables and values defining the state. Position and velocity
variables must be consistent with the selected coordinate system. The valid position
sets are given in \cref{tab:initial-position-variables} and the valid velocity sets in
\cref{tab:initial-velocity-variables}.

\begin{table}[htbp]
  \centering
  \caption{Initial position variables.}
  \label{tab:initial-position-variables}
  \begin{tabularx}{0.78\linewidth}{@{}l>{\ttfamily}lX@{}}
    \toprule
    Coordinate system & \normalfont Variable & Description \\
    \midrule
    Geodetic & alt  & Altitude \\
             & long & Longitude \\
             & lat  & Latitude \\
    \addlinespace
    Geocentric & rcm  & Distance from earth center \\
               & long & Longitude \\
               & lat  & Latitude \\
    \addlinespace
    ECFC or ECIC & x & Position $x$ component \\
                 & y & Position $y$ component \\
                 & z & Position $z$ component \\
    \bottomrule
  \end{tabularx}
\end{table}

Geodetic coordinates are the default. For example,

\begin{lstlisting}[style=taosmanual]
*initial
  alt=30000 long=0 lat=45       # Position vector variables
  vel=750 gamma=0 psi=90        # Velocity vector variables
  wt=2000 time=0                # Other variables
\end{lstlisting}

uses the geodetic position and velocity sets without explicitly naming the coordinate
system.

\begin{table}[htbp]
  \centering
  \caption{Initial velocity variables.}
  \label{tab:initial-velocity-variables}
  \begin{tabularx}{0.76\linewidth}{@{}l>{\ttfamily}lX@{}}
    \toprule
    Coordinate system & \normalfont Variable & Description \\
    \midrule
    Geodetic or geocentric & vel   & Velocity \\
                            & gamma & Flight-path angle \\
                            & psi   & Heading angle \\
    \addlinespace
    ECFC or ECIC & xdt & Velocity $x$ component \\
                 & ydt & Velocity $y$ component \\
                 & zdt & Velocity $z$ component \\
    \bottomrule
  \end{tabularx}
\end{table}

\begin{table}[htbp]
  \centering
  \caption{Other initial-condition variables.}
  \label{tab:other-initial-condition-variables}
  \begin{tabularx}{0.72\linewidth}{@{}>{\ttfamily}lX@{}}
    \toprule
    \normalfont Variable & Description \\
    \midrule
    mach    & Mach number, used instead of \statevar{vel} \\
    wt      & Weight \\
    mass    & Mass, used instead of \statevar{wt} \\
    time    & Time \\
    range   & Ground range \\
    path    & Path length \\
    t\_0     & Time at which \statevar{omega_0} applies \\
    omega\_0 & Angle between ECFC and ECIC at \statevar{t_0} \\
    \bottomrule
  \end{tabularx}
\end{table}

The units of all variables are those of the standard output variables in the Appendix
and may be changed with \problemblock{*units/fmt}. The only required value in
\cref{tab:other-initial-condition-variables} is \statevar{wt}, or equivalently
\statevar{mass}. One may replace the other, but both cannot be entered. Likewise,
\statevar{mach} may replace \statevar{vel} when specifying initial velocity, but
Mach number and velocity cannot both be supplied.

The remaining variables are optional and default to zero. The variables
\statevar{t_0} and \statevar{omega_0} control the angular difference between the
ECFC and ECIC coordinate systems. At time $t_0$, ECIC is rotated from ECFC by
$\Omega_0$; therefore, at mission time $t$,

\begin{equation}
  \Omega = \Omega_0 + \lVert\vec\omega_{\oplus}\rVert\,(t-t_0),
  \label{eq:initial-ecic-rotation}
\end{equation}

where the earth rotation rate is specified by \problemblock{*earth}.

The following example gives initial conditions in ECFC coordinates:

\begin{lstlisting}[style=taosmanual]
*initial ecfc
  x=19099345.2  y=9974312.7   z=6489105.4   # Position
  xdt=-173.973  ydt=10327.52  zdt=26425.9   # Velocity
  time=769.35   wt=1053.9                   # Time and weight
  t_epoch=-1045.2 omega_0=25.3              # ECFC to ECIC
\end{lstlisting}

The trajectory begins at 769.35~s. The ECFC-to-ECIC angle is not known at that
instant, so it is entered at the known time $-1045.2$~s, 1814.55~s before the
trajectory start. Times for all trajectories are referenced to the common $t=0$
mission time.

\taossubhead{Initial Conditions from a Trajectory}

A trajectory can obtain its initial conditions from another trajectory so that several
trajectories branch from a common vehicle. For example, if trajectory 1 is an aircraft
that releases a weapon at the beginning of segment 6, the weapon can be initialized
with

\begin{lstlisting}[style=taosmanual]
*initial from segment 6, trajectory 1
\end{lstlisting}

The \texttt{from} keyword requests initial conditions from another trajectory.
Conditions are always taken at the beginning of the specified segment.

In this form the \texttt{segment} and \texttt{trajectory} values are supplied as
shown, and position and velocity variables from
\cref{tab:initial-position-variables,tab:initial-velocity-variables} are omitted.
Values of \statevar{t_0} and \statevar{omega_0} may still follow the segment and
trajectory numbers and are often kept identical for every trajectory so all use the
same ECIC system.

After initialization from another trajectory, the first segment may use
\problemblock{*increment} and \problemblock{*reset} to adjust the copied state.
A deployed object's weight generally must be set because it differs from the original
vehicle, and its position is often adjusted because the two centers of mass are not at
exactly the same location.
